\properlane{interior-conversion}{The Soil of Confession, the Path of Grace, and the Altar Bridge}
\runninghead{The Soil of Confession and the Altar Bridge}

Between the historical Nativity already accomplished in Bethlehem and the public Judgment still to come at the end of time stands the present hour of the congregation gathered in church. On that middle ground, three textual questions press on the reader of the First Sunday of Advent. Why does the Alleluia ask God to \emph{show} His mercy (Ps.~84:8), as if mercy were invisible until revealed? How can sinful human soil yield \emph{its own} fruit (\latin{fructum suum}, Ps.~84:13) if every good comes from God? And why do the Introit and Gradual ask separately for God's ways (\latin{vias}) and His paths (\latin{semitas}, Ps.~24:4), while the Secret prays that these sacred Mysteries cleanse us to approach their own Source (\latin{principium})?

Six Catholic writers---five Fathers and Doctors together with Bl.~Ildefonso Schuster---answer those questions in terms of present conversion and sacramental purification. St.~Augustine expounds Psalms~24 and~84 as prevenient grace humbling pride and ploughing sinners' hearts through confession, and records in \work{Confessions}~VIII.12.29 how Romans~13:13--14 broke his own divided will; St.~Thomas Aquinas traces Psalm~24:1--4 from contemplative ascent through the common road of the commandments to the narrower footpath of the evangelical counsels; St.~Leo the Great preaches Advent fasting and almsgiving as cultivating the heart's soil; St.~Ambrose traces Christ's coming within each soul (\latin{in singulis}); St.~Bernard of Clairvaux calls this hidden visitation Christ's Middle Coming (\latin{adventus medius}); and Bl.~Ildefonso Schuster expounds the Secret, Communion, and Postcommunion of this Mass as the Eucharistic Mysteries purifying worshippers, refreshing arid hearts, and preparing the faithful for Christmas. Read together across the 1962 Mass, these expositions present Advent as a season of present interior conversion and Eucharistic renewal between Bethlehem and the Last Day.

\subsection{Prevenient Grace and the Plough of Confession}

St.~Augustine's exposition of Psalm~84 (\work{Enarrationes in Psalmos}~84.9, 14--16 [Coxe~85.6, 11--13]) takes a sharply different road from St.~Robert Bellarmine's Marian exegesis of the same psalm, solving both difficulties through the theology of prevenient and healing grace. On verse~8 (the Alleluia: ``Show us, O Lord, thy mercy; and grant us thy salvation''), Augustine explains why mercy must be \emph{shown}: fallen man naturally boasts as if his virtues were his own achievement. God shows His mercy when He persuades a man that whatever good he possesses comes wholly from God. Once a man sees that, Augustine traces an unbroken chain of spiritual stability: he is not proud; not being proud, he is not lifted up; not lifted up, he does not fall; not falling, he stands; standing, he clings fast; clinging fast, he abides and enjoys God (\S9 [Coxe~\S6]). Comparing the believer to a penniless wanderer told by a trustworthy messenger that he is a senator's son with a rich patrimony waiting at home, Augustine hears St.~John's promise (1~John~3:2) in \latin{salutare tuum da nobis}: ``grant us Thy Christ, that we may know Thy Christ, see Thy Christ; not as the Jews saw Him and crucified Him, but as the Angels see Him, and rejoice.''

When Augustine arrives at verses~12--14 (\S\S14--16 [Coxe~\S\S11--13]), he turns directly to the sinner's present transformation. Before God's righteousness can look down from heaven to pardon us, ``truth must spring out of the earth'' (v.~12). Sinful man is ``earth'' (Gen.~3:19): how can truth spring from an unrighteous man? ``Confess thy sins, and truth shall spring out of thee,'' Augustine answers. If a sinner calls himself righteous, he speaks falsehood; when he condemns his own sin truthfully, God's mercy replies: ``Let us spare this man, for he spareth not himself: let us pardon him, for he himself confesseth.''

At verse~13 (the Communion antiphon), Augustine's Old Latin Psalter reads \latin{Dominus dabit suavitatem} (``the Lord will give sweetness'') rather than the Gallican Vulgate's and Roman Missal's \latin{benignitatem} (``goodness''). That textual difference yields Augustine's explanation of how our earth bears fruit (\S15 [Coxe~\S12]). Sin holds people captive because it tastes sweet to fallen desire; mere commands cannot cure a divided heart until God pours into it a deeper interior sweetness in doing good. Then ``thou who at first didst delight in drunkenness, shall rejoice in sobriety: [\dots]\ and thou who didst take delight in robbing, shalt delight now in giving: thou whom shows delighted, shalt delight in prayer.'' Addressing his congregation directly, Augustine turns preaching and penance into farming: he has sown God's word in devout hearts ``furrowed, as it were, with the plough of confession,'' and now prays that God's rain will visit them at home, at meals, and in their walks so that ``our land shall give her increase.'' In verse~14, confession itself becomes the highway of John the Baptist (Matt.~3:3): before you confessed your sins, you had blocked God's road; confess your past life, and Christ has where to set His footsteps (\S16 [Coxe~\S13]). St.~Leo the Great (\work{Sermon}~17.1; cf.~\work{Sermon}~12.4) preaches the same Advent husbandry: fasting and almsgiving cultivate ``the ground of our heart'' so that it brings forth ``the germ of righteousness and the fruit of love.''

\subsection{Contemplative Ascent, Ways, and Paths}

St.~Augustine and St.~Thomas Aquinas approach Psalm~24:1--4 (the Introit, Gradual, and Offertory) in the same spirit of humble conversion. Asking who speaks in the psalm, Augustine (\work{Enarr.~in~Ps.}~24.1--4) answers with his doctrine of the whole Christ (\latin{totus Christus}): ``Christ speaks, but in the person of the Church,'' raising a soul once trodden down on the earth by carnal appetites. Why does the psalmist immediately pray \latin{non erubescam}, ``I shall not be ashamed''? Augustine traces shame back to Adam's fall: trusting in himself, man wished to be as God, and ended in bodily weakness, ``fearing death from the smallest insect,'' put to derision for his pride. Hidden spiritual enemies still whisper ``Well done, well done'' to lure believers off the narrow paths (\latin{semitae}) of Matthew~7:14, ``narrow, and known to few.''

St.~Thomas Aquinas (\work{Postilla super Psalmos}~24, Parma vol.~14, pp.~229--230) traces a clear progression through those four verses. Quoting St.~John Damascene, he defines lifting the soul (\latin{Ad te levavi}) as prayer itself---\latin{ascensus intellectus in Deum} (``the ascent of the intellect into God'') through contemplation of God's goodness and love---joined to unhesitating trust (\latin{in te confido}), since no one obtains what he asks without praying in faith (Jas.~1:6; Heb.~4:16). Aquinas distinguishes three causes of shame (\latin{erubescentia}): defeat by an adversary, frustrated hope (Rom.~5:5), and the memory of past sin (Rom.~6:21). Commenting on his Gallican text \latin{qui sustinent te}, he explains enduring in patient expectation by the very verb the Roman chant sings (\latin{qui cum longanimitate eum expectat}, citing Hab.~2:3, ``If it make any delay, wait for it: for it shall surely come''), and on the omitted verse~4a he draws from Ecclesiasticus~4:25 a distinction between ruinous confusion (eternal damnation) and salutary confusion (the shame of penance that brings glory and grace).

On verse~4, commenting on his Gallican text \latin{Vias tuas, Domine, demonstra mihi, et semitas tuas edoce me} (rather than the Gradual's Old Latin \latin{notas fac mihi}), Aquinas divides the petition inside the active life into two grades of obedience: \latin{viae} (``ways'') are the public, common highway of God's commandments (\latin{praecepta}, centered in charity, John~14:15; Matt.~19:17), whereas a \latin{semita} (``path'') is a narrower, shorter footpath (\latin{quaedam abbreviatio communis viae})---the evangelical counsels (\latin{consilia}, Jer.~6:16; Isa.~40:3). Because man does not know these ways on his own, the psalmist asks under \latin{Vias tuas} that God show them (\latin{ostende eas mihi}) both to the intellect and to the affections, before Aquinas turns under \latin{Et semitas} to the evangelical counsels and then to the contemplative guidance of verse~5 (\latin{Dirige me in veritate tua, et doce me}, ``Direct me in thy truth, and teach me''). St.~Robert Bellarmine (\work{Explanatio} on Ps.~24:4--5) underscores the same practical petition: the psalmist asks to be taught ``not speculatively, but practically; that is to say, he asks for such grace as may move his will to observe the commandments cheerfully.'' Read through Aquinas's distinction, hearing Psalm~24:4 twice before the Gospel---first with Gallican \latin{demonstra mihi} at the Introit and then with Old Latin \latin{notas fac mihi} at the Gradual---sets both the common road of the commandments and the narrower paths of the counsels before the assembly.

\subsection{Decisive Conversion and the Middle Coming at the Altar}

St.~Augustine's own conversion in \work{Confessions}~VIII.12.29 shows what happened when Romans~13:13--14 met a divided will. Torn between carnal habit and desire for God in the Milan garden, Augustine heard the command \latin{tolle, lege} (``take up and read''), opened the Apostle, and read the call to abandon rioting, drunkenness, chambering, and strife and ``put ye on the Lord Jesus Christ''---followed in his codex by verse~14b, ``and make not provision for the flesh, to fulfil the lusts thereof.'' Instantly, as the sentence ended, ``a light, as it were, of security'' flooded his heart and scattered every gloom of doubt. The 1962 Roman Missal ends at verse~14a, closing directly on Christ's name, whereas Augustine's conversion turned on reading through the end of the verse.

St.~Bernard of Clairvaux (\work{Sermones in Adventu}~3.3--4, pp.~25--26; 5.1--3, pp.~40--42) calls this hidden interior visitation the Middle Coming (\latin{adventus medius}). In Sermon~3.3--4 he distinguishes Christ's coming to men, His spiritual and invisible coming into men (citing John~14:23), and His coming against men; in Sermon~5.1 he places the hidden middle advent in spirit and virtue between the first coming in flesh and weakness (when Christ was our Redemption) and the third coming in glory and majesty (when He shall appear as our Life), calling it our Repose and Consolation---``the way, so to speak, by which we must travel from the first to the last.'' Bernard explains how to keep the Word so this coming takes hold (V.2--3): bread left in a cupboard can grow stale or be stolen by thieves or nibbled by mice, whereas bread eaten and digested passes into one's very strength; so the believer must receive God's Word into the mind, digest it into the affections and will, and let Christ rebuild the Jerusalem of the soul. St.~Ambrose (\work{Exp.~in~Luc.}~X.39--40) traces the same interior advent in the Gospel itself: Christ comes both at the world's end and \latin{in singulis}---within each soul that opens its affections to Him, shaking its powers with an increase of grace.

Bl.~Ildefonso Schuster (\work{The Sacramentary}~I, pp.~321--323) does not invoke Bernard's schema of a Middle Coming---on the Gospel (p.~322) he sets the first coming at Bethlehem directly beside the second coming in glory---yet his commentary on the prayers and Communion of this Mass points to the same present work of grace at the altar. In the Collect (p.~321), the Church asks our Lord to come to her aid and save her by His power from the dangers brought on by sin (\latin{imminentibus peccatorum nostrorum periculis}), depending wholly on His protection and deliverance. At the Offertory (\latin{Ad te levavi}, p.~322), the assembly waits upon the Lord in watchful trust so as never to be confounded. In the Secret (\latin{Haec sacra nos}, p.~322), Schuster hears ``the ancient idea that runs through the Roman epiklesis,'' beseeching the Savior that the power of the Eucharistic sacrifice purify our minds so we may worthily receive Him who is its Author (\latin{principium}). At Communion (Ps.~84:13, pp.~322--323), Schuster reads the Eucharist as ``a pledge of his infinite goodness'' refreshing hearts rendered arid and sterile by sin with the dews of grace so that they bring forth fruit in due season (a moral application echoed in William Durandus, \work{Rationale}~VI.3, fol.~261v, where \latin{terra nostra} is our flesh yielding good works). In the Postcommunion (\latin{Suscipiamus}, p.~323), which Schuster calls ``the true \latin{Eucharistia} in its etymological sense,'' he writes that the divine virtue of the heavenly bread helps worshippers prepare suitably for the approaching feast of our common ransom. Read alongside Bernard's Middle Coming, these prayers show the sacramental path by which Christ purifies the assembly between Bethlehem and Judgment: vested in violet without \latin{Gloria in excelsis}, the Church worships the Triune God in the Sunday Preface and is sent back into daily life at \latin{Ite, missa est} (RGMR~507).

\subsection{Tensions and Limits of This Perspective}

Two textual facts qualify how St.~Augustine's words bear on the 1962 Mass. First, Augustine's account of interior ``sweetness'' (\latin{suavitas}) transforming sinful desires in Psalm~84:13 rests on his Old Latin Psalter, whereas the Roman Missal sings the Gallican word \latin{benignitatem} (``goodness''); the theology of grace healing human delight accords with the verse, yet the verbal hook differs from the Missal's printed chant. Second, Augustine's garden conversion in \work{Confessions}~VIII.12.29 turned directly on Romans~13:14b, which the 1962 Epistle omits. Finally, interpreting the Gospel primarily through Ambrose's \latin{in singulis} (Christ coming within each soul) would narrow the cosmic, public scale of Luke~21:25--27 if detached from the universal judgment expounded by Gregory and Cyril.

\subsection{The Four Senses}

\begin{fourSenses}
\item[Literal.] A believer prays for deliverance from shame and guidance in God's law (Ps.~24); restored Israel cries for God's mercy and fruitful land (Ps.~84); St.~Paul exhorts Roman Christians to leave behind carnal excess and put on Christ (Rom.~13:11--14a); Jesus urges fearless constancy grounded in His unfailing words (Luke~21:25--33); the Roman Church begs for cleansing through the Mysteries and reverent preparation for the coming feast (see \hyperref[sec:date-location]{\work{Scriptural Date and Location}} for scriptural chronology).
\item[Allegorical.] Christ speaks in the person of His Body the Church (\latin{totus Christus}, Augustine on Ps.~24); He visits souls invisibly within each believer by an increase of grace (Ambrose~X.39--40) and in His hidden Middle Coming in spirit and virtue (Bernard, \work{In~Adventu}~3.3--4, 5.1--3); and the Eucharistic Mysteries act as a purifying epiklesis and a pledge of divine goodness refreshing arid hearts (Schuster on the Secret and Ps.~84:13, pp.~322--323).
\item[Moral.] Prevenient grace curing pride so man credits all good to God (Augustine, Ps.~84:8); confession ploughing hardened hearts so interior delight in righteousness displaces delight in sin (Augustine, Ps.~84:13); decisive abandonment of entrenched vice to put on Christ (Augustine, \work{Confessions}~VIII.12.29); progressing from the common road of commandments (\latin{viae}) to the narrower paths of the evangelical counsels (\latin{semitae}, Aquinas, Ps.~24:4); Advent fasting, prayer, and almsgiving cultivating the heart's ground (Leo, Serm.~12.4, 17.1).
\item[Anagogical.] Christ's hidden Middle Coming in grace and the Eucharist serving as the bridge that purifies worshippers now so they may stand unashamed (\latin{non confundentur}) before their divine \latin{principium} at His third coming in glory.
\end{fourSenses}
